%% kapa: exam=1 round=25 year=2014 period=2 form=A subject=accounting subject_ko=회계학 seq=1 number=41
%% source: Q-Net 감정평가사 기출문제 (cst00309) artlSeq=5101207 file=extracted/1008/5101207/2014년 제25회 감정평가사 1차시험 기출문제/2교시(A형).hwp
%% conversion: kordoc 4.24.51 -> tools/gen1.py (auto); review=proofread-sonnet; flags=none
\begin{question}{41}
재무보고를 위한 개념체계상 유용한 정보의 질적 특성에 관한 설명으로 옳지 않은 것은?
\begin{choices}
  \item 충실한 표현은 모든 면에서 정확한 것을 의미하지는 않는다.
  \item 중요성은 개별 기업 재무보고서 관점에서 해당 정보와 관련된 항목의 성격이나 규모 또는 이 둘 모두에 근거하여 해당 기업에 특유한 측면의 목적적합성을 의미한다.
  \item 완전한 서술은 필요한 기술과 설명을 포함하여 정보이용자가 서술되는 현상을 이해하는 데 필요한 모든 정보를 포함하는 것이다.
  \item 중립적 정보는 목적이 없거나 행동에 대한 영향력이 없는 정보를 의미하지 않는다.
  \item 재무정보의 비교가능성은 정보이용자가 항목 간의 차이점을 식별하고 이해할 수 있게 하는 질적 특성으로 비슷한 것을 달리 보이게 함으로써 보강된다.
\end{choices}
\end{question}
\answer{5}
